%%
%% fall-lecture08.tex
%% 
%% Made by Alex Nelson <pqnelson@gmail.com>
%% Login   <alex@lisp>
%% 
%% Started on  2026-10-15T10:17:56-0700
%% Last update 2026-10-15T10:17:56-0700
%% 

\lecture[Complete Metric Spaces]{}

\begin{node}
The second problem set: Folland Chapter~4 problems 37, 43, 49 [Let $X$
  be a compact Hausdorff space and $E\subset X$; (a) If $E$ is open,
  then $E$ is locally compact in the relative topology; If $E$ is
  dense in $X$ and locally compact in relative topology, then $E$ is
  open; (c) $E$ is locally compact in relative topology iff $E$ is
  relatively open in $\closure{E}$], 52
[prove the one-point compactification of $\RR^{n}$ is homeomorphic to
  the $n$-sphere], 77 [countable product of metric spaces is a metric space].
\end{node}

\begin{definition}
Let $(X,\rho)$ be a metric space (with metric topology). We say it is
\define{Complete} if for any Cauchy sequence $(x_{n})$ [i.e., any
  sequence such that $\forall m,n\in\NN\ldotp\rho(x_{m},x_{n})\to0$ as $m,n\to\infty$]
has a limit $x_{n}\to x$.
\end{definition}

\begin{example}
The real line $\RR$ with the standard topology is complete. But the
subspace $\QQ$ is not complete.
\end{example}

\begin{example}
The circle $S^{1}$ is complete, but when we remove a point
$S^{1}\setminus\{\mbox{pt}\}$ it is not complete. So what? Well,
$S^{1}\setminus\{\mbox{pt}\}$ is homeomorphic to $\RR$ with the
standard topology (which \emph{is} complete). This tells us that
``completeness'' \emph{is not a topological property}.
\end{example}

\begin{theorem}[Banach's fixedpoint]
Let $(X,\rho)$ be a complete metric space, let $T\colon X\to X$ be a
strict contraction [there exists some $\alpha\in(0,1)$ such that for
  all $x,y\in X$ we have $\rho(T(x),T(y))<\alpha\rho(x,y)$].
Then there exists a unique fixedpoint $x_{*}\in X$ such that $T(x_{*})=x_{*}$.
\end{theorem}

[Observe the orbit of a point is a Cauchy sequence in $X$.]

This is useful when it comes to studying differential equations.

\begin{definition}
Let $(X,\rho)$ be a metric space.
Then the \define{Completion} of $X$ is $(\widetilde{X},\widetilde{\rho})$
where $\widetilde{X}$ is the set of all equivalence classes of Cauchy
sequences where $(x_{n})\sim(y_{n})$ iff $\rho(x_{n},y_{n})\to0$ (so
eventually, there is some $N$ such that $x_{n}=y_{n}$ for all $n\geq N$). We
define the metric $\widetilde{\rho}$ on $\widetilde{X}$ by
\begin{equation}
\widetilde{\rho}([(x_{n})],[(y_{n})]):=\lim_{n\to\infty}\rho(x_{n},y_{n}),
\end{equation}
where $[(x_{n})]$, $[(y_{n})]$ are the equivalence classes of Cauchy
sequences with representatives $(x_{n})$ and $(y_{n})$ respectively.

We also have the isometry $X\into\widetilde{X}$ sending $x\mapsto[(x_{n}=x)]$
points to the constant sequence.

We need to prove this is all well-defined (especially: the limit in
the definition of $\widetilde{\rho}$ exists, and that the image of
$X\into\widetilde{X}$ is dense in $\widetilde{X}$). We claim that
$(\widetilde{X},\widetilde{\rho})$ is a complete metric space.
\end{definition}

\subsection*{Compactness of metric space}

\begin{node}
Recall in $\RR^{n}$, we have $X\subset\RR^{n}$ is compact if and only
if $X$ is closed and bounded [Heine--Borel theorem]. The extension of
this statement to general metric spaces is the following theorem.
\end{node}

\begin{theorem}[Hausdorff]
Let $(X,\rho)$ be a metric space. Then the following are equivalent:
\begin{enumerate}
\item $X$ is compact
\item $X$ is sequentially compact
\item $X$ is complete and ``totally bounded'' (see next definition).
\end{enumerate}
In particular, for metric spaces, ``compactness'' coincides with
``sequential compactness''.
\end{theorem}

\begin{definition}
Let $(X,\rho)$ be a metric space. We say $X$ is \define{Totally Bounded}
if for each $\varepsilon>0$ there exists a finite cover of $X$ by open
balls of radius $\varepsilon$.
\end{definition}

\begin{proof}[Hausdorff theorem]
$(1)\implies(2)$ Assume $X$ is compact.
Let $(x_{n})$ be a sequence in $X$.
[We want to find a convergent subsequence.]
Assume for contradiction there is no convergent subsequence, then
$(x_{n})$ has no limit points. So for every $x\in X$, there is an open
neighborhood $U_{x}\ni x$ such that $x_{n}\notin U_{x}$ for $n>N_{x}$ (for
some $N_{x}\in\NN$ which depends on $x$). Then as we vary $x$, we get
an open cover $\{U_{x}\}_{x\in X}$ and by compactness there is a
finite subcover $U_{x_{1}}$, \dots, $U_{x_{n}}$. Then for large $N\geq\max(N_{x_{1}},\dots,N_{x_{n}})$,
we have $x_{m}\notin U_{x_{i}}$ for all $i=1,\dots,n$. So $x_{m}\notin U_{x_{1}}\cup(\cdots)\cup U_{x_{n}}$
which is equivalent to $x_{m}\notin X$, which is our desired contradiction.

$(2)\implies(3)$ Assume $(X,\rho)$ is sequentially compact.

\textsc{Claim 2.1: $X$ is complete}. Let $(x_{n})$ be a Cauchy
sequence in $X$. Then by sequential compactness, there is a convergent
subsequence $x_{n_{k}}\to x$. Then by the triangle inequality and
Cauchy property, we see that $x_{n}\to x$.

\textsc{Claim 2.2: $X$ is totally bounded}.
Assume for contradiction that $X$ is not totally bounded (but still
sequentially compact). Then there exists an $\varepsilon>0$ such that
there are infinitely many points which are a distance greater than
$\varepsilon$ from each other. (This is by induction: take $x_{1}$,
then a point $x_{2}$ such that $\rho(x_{1},x_{2})>\varepsilon$; then
take $x_{3}$ such that $\rho(x_{1},x_{3})>\varepsilon$ and
$\rho(x_{2},x_{3}>\varepsilon$; \dots) Then this $X$ cannot be sequentially
compact, because we just constructed a sequence $(x_{n})$ with no
convergent subsequence. Hence our contradiction.

$(3)\implies(1)$ Assume $(X,\rho)$ is both complete and totally bounded.
(Note: total boundedness implies separability since total boundedness
implies the existence of a countable dense set.) Then recall
separability is equivalent to second-countability for metric spaces,
and second-countability is equivalent to Lindel\"{o}f for metric
spaces. This means that we have for any open cover of $X$ a countable open subcover for $X$.
Great, enough recalling.

We want to prove every open cover has a finite subcover. Let
$X=\bigcup_{n}\mathcal{G}_{n}$ be a countable open cover of $X$.
Without loss of generality, we can assume the sequence of open sets
$\mathcal{G}_{n}$ is increasing [if not, just define
  $U_{n}=\bigcup^{n}_{k=1}\mathcal{G}_{k}$ and use the $U_{n}$ as the
  open cover which is increasing].
We want to prove after finitely many steps, $\mathcal{G}_{N}=X$ for
some $N\in\NN$.

Consider instead the [decreasing] sequence of $F_{n}=X\setminus\mathcal{G}_{n}$
closed sets. Then $F_{n}$ is monotone decreasing and
\begin{equation}
\bigcap_{n}F_{n}=\emptyset.
\end{equation}
We want to find an $N\in\NN$ such that $F_{N}=\emptyset$.

Assume for contradiction there is no $N$ such that $F_{N}=\emptyset$.
We can prove that $X$ is sequentially compact (which we defer to the end of the proof).
If for all $n$ we have $F_{n}\neq\emptyset$, then take $x_{n}\in F_{n}$.
Then $x_{n+1}\in F_{n+1}\subset F_{n}$,
and $x_{n+2}\in F_{n+2}\subset F_{n+1}\subset F_{n}$. So this gives us
a sequence $(x_{n})$ such that $\forall n\in\NN\ldotp x_{n}\in F_{n}$.
Then by sequential compactness, there is a subsequence $x_{n_{k}}\to x$
where $x\in\bigcap_{n}F_{n}$. But this means $\bigcap F_{n}\neq\emptyset$.
But this contradicts the property $\bigcap F_{n}=\emptyset$. Hence we
reject our assumption and conclude there is an $N\in\NN$ such that
$F_{N}=\emptyset$. 

\textbf{Deferred proof} ($X$ is sequentially compact)
Let $(x_{n})$ be a sequence in $X$. Take $\varepsilon=1/2$. Then by
total boundedness, $X$ can be covered by finitely-many radius-$1/2$ open balls
and (by pidgeon hole principle, at least) one of these open balls must contain
infinitely many elements of the sequence.
Take one of these open balls $B(r=1/2, x)$ containing infinitely many
points of the sequence. We then iterate, with $\varepsilon=2^{-2}$,
cover $X$ be radius-$2^{-2}$ balls, then looking at the balls covering
$B(r=2^{-1},x)$ we see there are finitely many such radius-$2^{-2}$
balls and one of them must contain infinitely many terms in the
sequence. We can iterate as needed, and this gives us a way to
construct a subsequence $x_{n_{k}}$ of $x_{n}$ which is Cauchy (hence,
by completeness of $X$, converges to some
point of $X$). This proves subsequential compactness of $X$.
\end{proof}